%Dokumentklasse
\documentclass[a4paper,12pt]{scrreprt}
\usepackage[left= 2.5cm,right = 2cm, bottom = 4 cm]{geometry}
 \usepackage[onehalfspacing]{setspace}
% ============= Packages =============

% Dokumentinformationen
\usepackage[
	pdftitle={Testautomatisierung},
	pdfsubject={},
	pdfauthor={Matthias Karl},
	pdfkeywords={}
	pdftex=true, 
	colorlinks=true,
 	breaklinks=true,
	citecolor=black,
	linkcolor=black,	
	menucolor=black,	
	urlcolor=black
]{hyperref}


% Standard Packages
\usepackage[utf8]{inputenc}
\usepackage[ngerman]{babel}
\usepackage[T1]{fontenc}
\usepackage{graphicx}
\graphicspath{{img/}}
\usepackage{fancyhdr}
\usepackage{lmodern}
\usepackage{color}
\usepackage{transparent}

% zusätzliche Schriftzeichen der American Mathematical Society
\usepackage{amsfonts}
\usepackage{amsmath}

% nicht einrücken nach Absatz
\setlength{\parindent}{0pt}


% ============= Kopf- und Fußzeile =============
\pagestyle{fancy}
%
\lhead{}
\chead{}
\rhead{\slshape \leftmark}
%%
\lfoot{}
\cfoot{}
\rfoot{\thepage}
%%
\renewcommand{\headrulewidth}{0.4pt}
\renewcommand{\footrulewidth}{0pt}

% ============= Package Einstellungen & Sonstiges ============= 

%Besondere Trennungen
\hyphenation{De-zi-mal-tren-nung St-rei-fen-licht-scan-nern}

%römische Aufzählungen mit \RM{Zahl}
\newcommand{\RM}[1]{\MakeUppercase{\romannumeral #1}}


% ============= Dokumentbeginn =============

\begin{document}

\pagestyle{empty}
\begin{center}
\begin{tabular}{p{\textwidth}}


\begin{center}
\includegraphics[scale=0.17]{img/logos.jpg}
\end{center}


\\

\begin{center}
\LARGE{\textsc{
Testautomatisierung \\  \large{- mit Selenium -}
}}
\end{center}

\\


\begin{center}
\large{Fakultät für Informatik und Mathematik \\
der Hochschule München \\}
\end{center}

\\

\begin{center}
\textbf{\Large{Seminararbeit}}
\end{center}



\begin{center}
vorgelegt von
\end{center}

\begin{center}
\large{\textbf{Matthias Karl}} \\
\large{Matrikel-Nr: 03280712} \\
\end{center}

\begin{center}
\large{im Dezember 2014}
\end{center}

\\

\\

\begin{center}
\begin{tabular}{lll}
\textbf{Prüferin:} & & Prof. Dr. Ulrike Hammerschall\\
\end{tabular}
\end{center}

\end{tabular}
\end{center}

% \part im Inhaltsverzeichnis nicht nummerieren
\makeatletter
\let\partbackup\l@part
\renewcommand*\l@part[2]{\partbackup{#1}{}}

%Seitennummerierung neu beginnen, Zahlen [arabic], röm.Zahlen [roman,Roman], Buchstaben [alph,Alph]
\pagenumbering{Roman}

\newpage
\addsec{Eidesstattliche Erklärung}
\label{erklaerung}

Hiermit versichere ich, die vorliegende Studienarbeit selbstständig und nur unter Verwendung der von mir angegebenen Quellen und Hilfsmittel verfasst zu haben. Sowohl inhaltlich als auch wörtlich entnommene Inhalte wurden als solche kenntlich gemacht. Die Arbeit hat in dieser oder vergleichbarer Form noch keinem anderem Prüfungsgremium vorgelegen. \\
\\[1.5cm]
Datum:	\hrulefill\enspace Unterschrift: \hrulefill
\\[3.5cm]
 



\addsec{Zusammenfassung / Abstract}
\label{sec:zusammenfassung}
Diese Arbeit befasst sich mit der Automatisierung von Testfällen. Die Arbeit gibt einen Überblick über die verschiedenen Bereiche und Möglichkeiten bei der Testautomatisierung und befasst sich im nähren mit dem Automatisierungstool Selenium.
Zu Beginn werden zunächst die allgemeinen Grundlagen des Testens geklärt um
dann die verschiedenen Bereiche und Möglichkeiten der Testautomatisierung aufzuzeigen. Auf den Bereich der automatisierten UI-Tests wird im Kontext von Selenium zum Schluss näher eingegangen.



\newpage
\pagestyle{fancy}
%Inhaltsverzeichnis
\tableofcontents

\newpage
%Seitennummerierung neu beginnen, Zahlen [arabic], röm.Zahlen [roman,Roman], Buchstaben [alph,Alph]
\pagenumbering{arabic}
% pagestyle für gesamtes Dokument aktivieren
\pagestyle{fancy}

\newpage
\chapter{Einleitung}
\label{sec:einleitung}

Software hat in der heutigen Zeit eine hohe Verbreitung gefunden. Softwaresysteme werden immer wichtiger. Sowohl für Unternehmen als auch für jeden einzelnen persönlich. Die Anforderungen an moderne Software steigen ständig. Die Systeme werden immer größer und komplexer. Das hat zur Folge, dass auch die Anforderungen an die Qualität der Software immer weiter Steigen und wichtiger werden.
Fehler in Software verursachen immer wieder einen hohen finanziellen Schaden und können im schlimmsten Fall sogar Menschenleben kosten. Diese Probleme werden immer gravierender wenn mit der Wichtigkeit, Komplexität und Größe von Software nicht auch gleichzeitig die Qualität steigt. \cite{burnstein_practical_2003} \newline
Softwaretests sind ein weit verbreitetes Mittel um die Qualität einer Software zu überprüfen und sicherzustellen.
Komplexere und größere Software bedeutet daher auch gleichzeitig einen steigenden Testaufwand.\newline
Qualitätssicherungsmaßnahmen wie beispielsweise das Testen machen jetzt schon einen Großteil der Kosten in Softwareprojekten aus. Studien haben gezeigt, dass das Testen für 50\% und mehr der gesamten Projektkosten verantwortlich ist. \cite{ramler_economic_2006} 
Es ist daher nicht verwunderlich, dass es immer wieder versuche gibt diese Kosten zu reduzieren ohne dabei den angestrebten Qualitätsstandart zu reduzieren.
Einer der Wege die dafür vorgeschlagen wurden ist das automatisieren von Testfällen 
\cite{harrold_testing:_2000}
Im laufe der Jahre hat die Teatautomatisierung immer mehr an Bedeutung gewonnen. Heute ist sie bereits fester Bestandteil von Bereichen wie Continous Delivery und Continous Integration.

Diese Arbeit befasst sich daher mit der Automatisierung von Testfällen. Die Arbeit soll einen Überblick über die verschiedenen Bereiche der Testautomatisierung geben um dann den Bereich der Automatisierten GUI-Tests näher zu beleuchten.
Kapitel 1 ...







\include{05_architekturvarianten_des_cloud_computing}

\include{06_anwendungen_von_docker}

\include{07_fazit}


%Verzeichnis aller Bilder
\newpage
\listoffigures

%Literaturverzeichnis
\newpage
\bibliographystyle{alphadin}
\bibliography{Literatur}

\end{document}
